This thesis studies the problem of robust strategy optimization in systematic trading. The central
question is how trading strategies can be optimized on past market data while limiting the risk of
overfitting, but more specifically, which optimization procedure produces more reliable
out-of-sample performance.

To address this question, the project develops an efficient research framework composed of a C++
backtesting and optimization engine linked to Python reporting and analysis tools. The C++ component
is responsible for data processing, strategy simulation, trade reconstruction, and grid-based
parameter optimization. The Python component is used for plotting, report generation, and
performance analysis. This separation allows the computationally intensive parts of the workflow to
remain efficient and independent, while keeping the analytical layer flexible.

The empirical analysis uses the same discretionary optimization procedure under two different
data-generation settings. The benchmark approach optimizes the strategy on the historical in-sample
period, defined as the first 70\% of the original series. The alternative approach optimizes the
strategy on synthetic markets generated from the same in-sample window: multiple synthetic paths are
produced, the strategy is evaluated across them, and parameter selection is based on average
performance rather than on a single historical path. In both cases, parameter selection is performed
using discretion, following basic robustness principles rather than blindly selecting the single
best result. The selected parameters are then frozen and tested on the same out-of-sample period,
corresponding to the final 30\% of the original historical series. The comparison focuses on the gap
between optimization performance and out-of-sample performance: a smaller deterioration is
interpreted as evidence of lower overfitting.

The results do not show a clear reduction in overfitting from the synthetic optimization method. The
synthetic paths are useful as alternative market scenarios, but in these experiments they do not lead
to parameters that generalize better than the standard IS/OOS benchmark.
